\documentclass[11pt,a4paper]{article}
\usepackage[T1]{fontenc}
\usepackage[utf8]{inputenc}
\usepackage{lmodern}
\usepackage[margin=27mm,headheight=15pt]{geometry}
\usepackage{amsmath,amssymb,amsthm,mathtools}
\usepackage{microtype,booktabs,longtable,array,enumitem}
\usepackage{xcolor,listings,fancyhdr,xurl}
\usepackage[colorlinks=true,linkcolor=blue!45!black,citecolor=blue!45!black,urlcolor=blue!45!black]{hyperref}
\definecolor{ink}{HTML}{17324D}
\hypersetup{pdftitle={A uniform introduction rule for intuitionistic primitive recursive arithmetic},pdfauthor={Research prepared for Vladimir Reshetnikov},pdfsubject={An explicit seven-rule calculus, bi-interpretability, and whole-checker comparisons}}
\pagestyle{fancy}
\fancyhf{}
\fancyhead[L]{\small\textsc{Uniform introduction for iPRA}}
\fancyhead[R]{\small November 11, 2026}
\fancyfoot[C]{\thepage}
\setlength{\parskip}{4pt}
\setlength{\parindent}{0pt}
\setlength{\emergencystretch}{2em}
\setlist{itemsep=3pt,topsep=5pt}
\newtheorem{theorem}{Theorem}[section]
\newtheorem{lemma}[theorem]{Lemma}
\newtheorem{proposition}[theorem]{Proposition}
\newtheorem{corollary}[theorem]{Corollary}
\theoremstyle{definition}
\newtheorem{definition}[theorem]{Definition}
\newtheorem{remark}[theorem]{Remark}
\newcommand{\ipra}{\mathrm{iPRA}}
\newcommand{\Teq}{T_{\mathrm{eq}}}
\newcommand{\FV}{\operatorname{FV}}
\newcommand{\eqv}{\dashv\vdash}
\lstset{basicstyle=\ttfamily\footnotesize,columns=fullflexible,keepspaces=true,breaklines=true,frame=single,framerule=.3pt,rulecolor=\color{black!20},showstringspaces=false,aboveskip=10pt,belowskip=10pt}
\begin{document}
\begin{titlepage}
\vspace*{12mm}
{\color{ink}\Large\textsc{Concise iPRA research}}\par
\vspace{9mm}
{\color{ink}\huge\bfseries A uniform introduction rule for intuitionistic primitive recursive arithmetic\par}
\vspace{7mm}
{\Large An unrestricted seven-rule calculus,\par both interpretation composites,\par and a measured checker comparison\par}
\vspace{12mm}
{\large Prepared for Vladimir Reshetnikov\par November 11, 2026\par}
\vfill
\textbf{Result and limitation.}
The new calculus has a single introduction schema whose mode depends only
on the matrix root. Every equation, $Q$ formula, and $D$ formula belongs to
its public language. A reversed iterator orbit supplies the missing
progress elimination, and the report proves bi-interpretability with
full first-order iPRA.

This is a structural simplification, not a proof of a globally smallest
axiomatization. In matched straightforward implementations it is smaller
than the adopted eight-rule kernel, but larger than the already known
seven-rule witness-sensitive kernel. A separate, small implementation
reduction improves the inspected compact envelope; it is not an
implementation of the new orbit calculus.\par
\vspace{8mm}
\textbf{Reference snapshot.}\quad
\texttt{VladimirReshetnikov/ProveIt}\par
\texttt{a4e5e96c456303cf231eab2aeeb3291d8889833f}\par
\path{Logic/PrimitiveRecursiveArithmetic/ConciseIPRA}
\end{titlepage}

\begin{abstract}
We replace the separate case and equation-induction introductions of the
current native scalar kernel by one introduction. At an equation matrix,
$Q_y E$ means $\forall y(E(y)\to E(Sy))$; at a compound matrix it keeps the
conditional reading $B(0)\to\forall y B(y)$. The corresponding $D$ binds an
existential witness. Elimination specializes a closure along
$I(a,n;u.Su)$, with the source placed in the counter position. Its zero
source computes directly to the arbitrary target $n$, which derives
contextual equation induction without a circular use of an iterator
identity. The new rules admit numerical substitution by direct induction.
Explicit translations preserve derivations in both directions and have
provably identity formula composites for the entire language, not only
guarded public formulas. Numerical domains remain whole and equality
remains actual equality. The written theorem is accompanied by a standalone
typed proof checker, matched old kernels, reproducible size measurements,
and positive and negative certificates. The new metatheorems are not
claimed to have been checked in Lean or Rocq.
\end{abstract}

{\small\tableofcontents}
\newpage
% BEGIN sections/context.tex
\section{The comparison target and the meaning of simplicity}

The reference is the exact repository snapshot displayed on the title
page, inspected through GitHub. Its main native kernel $\mathcal N$
has equations, a one-variable closure $Q$, and a two-variable binder $D$.
It has eight primitive rule constructors: hypothesis, two computation
equations, equational transport, $Q$ elimination, $D$ elimination,
$D$ introduction by cases, and $Q$ introduction by equation progress.
The repository proves its equivalence with full first-order iPRA in Lean
\cite{baseline,native,bridge}.

The reference also contains a seven-rule presentation $W$.
Its single introduction uses progress exactly when the matrix is an
equation and the witness is the literal zero; it uses cases otherwise.
The theorem \texttt{wprf\_iff\_nprf} proves that this presentation has the
same sequents as the native kernel. Consequently, a claim of novelty
based on the number seven alone would be incorrect. A numerical
substitution can change a witness into literal zero, so a direct induction
on $W$ derivations does not preserve the primitive introduction choice
\cite{witness}.

The present construction changes the meaning and elimination of atomic
closures. This makes the introduction choice depend only on the formula
root, which numerical substitution preserves. Unlike the guarded
presentations \cite{guarded}, every formula generated by the displayed
grammar is public. No restriction to a prescribed zero guard, protected
equation, or mirror shape is imposed.

\begin{center}
\small
\begin{tabular}{@{}p{.18\textwidth}p{.18\textwidth}p{.29\textwidth}p{.24\textwidth}@{}}
\toprule
Presentation & Introductions & Primitive mode selection & Public formulas \\
\midrule
Adopted $\mathcal N$ & Cases and equation progress & Two rules & All Eq/$Q$/$D$ \\
Known $W$ & One & Matrix root and literal witness & All Eq/$Q$/$D$ \\
Guarded core & One & Matrix root & Restricted atomic matrices \\
New $\mathcal R$ & One & Matrix root only & All Eq/$Q$/$D$ \\
\bottomrule
\end{tabular}
\end{center}

The scalar terms, arithmetic construction, and use of progressive
conditions are inherited ideas. In particular, an earlier progressive
presentation uses progress at every matrix root, counter-aware recursion,
separate step and iteration eliminations, and an explicit Boolean-cases
rule \cite{progressive}. Here the iterator remains scalar, compound
closures retain their conditional reading, there is no primitive
Boolean-cases rule, and one orbit specialization serves both kinds of
closure. The all-formula normalization theorem is the substantive new
result.

\subsection{What the size claim is}

The relevant implementation criterion is the size of a complete
standalone proof checker, with gzip level 9 reported separately for
information. We measure significant Python tokens, all AST nodes, and
UTF-8 bytes after removing genuine docstrings and applying
\texttt{ast.unparse}, followed by one newline. No producer or theorem
database is imported by a checker.

Two comparisons are needed. The primary comparison uses explicit tagged
syntax, exact integer/list checks, de Bruijn scoping, and complete
Eq/$Q$/$D$ statement validation in three matched standalone programs.
The secondary comparison concerns the repository's much smaller arity
codec under its canonical-public-goal convention. That convention does
not include a complete public-grammar validator. Its numbers cannot be
interchanged with those of the typed interface.

The inspected incoming archives are part of the comparison, not omitted
competitors. Their original ZIP bytes were authenticated against the six
Git blob identifiers at the pinned commit. They include
\texttt{deterministic\_auxiliary}, \texttt{guarded\_source},
\texttt{guarded\_orbit}, \texttt{mirror\_guard},
\texttt{protected\_core}, and \texttt{reversal\_core}
\cite{incoming}. Exact source records and the comparison measurements
are included in the archive.

\textbf{Scope of the outcome.}
The new theory gives a cleaner introduction discipline and preserves full
bi-interpretability. It does not improve every existing checker-size
coordinate. The separate compact implementation gain is small and does
not establish a new mathematical lower bound. No global minimality
theorem is claimed.

% END INCLUDE

% BEGIN sections/new_theory.tex
\providecommand{\llbracket}{[\![}
\providecommand{\rrbracket}{]\!]}
\section{A seven-rule calculus with complete progressive elimination}
\label{sec:new-calculus}

The new presentation keeps the scalar iterator and the three native formula
constructors.  It changes the meaning of an atomic closure: an atomic closure
records a progressive equation, while a compound closure retains the old
conditional-universal meaning.  A single elimination rule recovers both
one-step progress and induction from an atomic closure.  Consequently every
formula of the new language has a faithful ordinary reading; no restricted
public grammar or auxiliary-only formula class is needed.

This is a structural simplification, not a claim of an absolute minimum.  The
current Lean kernel has eight rule constructors.  The proposed kernel has
seven, and its introduction rule commutes directly with numerical
substitution.  The earlier seven-rule \texttt{WPrf} presentation also has seven
constructors, but its introduction distinguishes the literal witness $0$ from
other witnesses.  Here the distinction depends only on whether a matrix is a
literal equation, which numerical substitution cannot change.  The concrete
checker measurements, including comparisons in which this calculus is larger,
are reported separately.\cite{native,witness}

\subsection{Complete primitive specification}
\label{subsec:new-specification}

There is one numerical sort.  Terms are
\[
 t ::= x\mid 0\mid St\mid I(n,a;u.h).
\]
Only $u$ is bound in $h$; other occurrences are free parameters.  Substitution
is capture-free.  Put
\begin{equation}
 1:=S0,\qquad
 \mathsf A(a,n):=I(a,n;u.Su).
 \label{eq:new-add}
\end{equation}
Thus the first argument of $\mathsf A$ is the iteration counter and the
second is its initializer.  This argument order is essential below.

Formulas are
\begin{equation}
 F ::= (s=t)\mid Q_yF\mid D_{x,y}F.
 \label{eq:new-formulas}
\end{equation}
The binder $Q$ binds $y$, and $D$ binds the distinct variables $x,y$.
A \emph{literal equation} has outermost constructor $=$.  A \emph{compound}
formula has outermost constructor $Q$ or $D$.  Define the syntactic selector
\begin{equation}
 \eta_B(a):=
 \begin{cases}
 a,&B\text{ is a literal equation},\\
 0,&B\text{ is compound}.
 \end{cases}
 \label{eq:new-eta}
\end{equation}
Write $\Gamma\vdash_{\mathcal R}F$ for derivability.  Contexts are finite
reusable lists of formulas.  The complete primitive inventory is the following
seven rule schemas:
\[
 \frac{F\in\Gamma}{\Gamma\vdash_{\mathcal R}F}
 \quad\mathrm{Hyp},
\]
\[
 \frac{}{\Gamma\vdash_{\mathcal R}I(0,a;u.h)=a}
 \quad\mathrm{I0},\qquad
 \frac{}{\Gamma\vdash_{\mathcal R}
 I(Sn,a;u.h)=h[I(n,a;u.h)/u]}
 \quad\mathrm{IS},
\]
\begin{equation}
 \frac{\Gamma\vdash_{\mathcal R}s=t\qquad
       \Gamma\vdash_{\mathcal R}p(s)=q(s)}
      {\Gamma\vdash_{\mathcal R}p(t)=q(t)}
 \quad\mathrm{Tr}_{=},
 \label{eq:new-tr}
\end{equation}
\begin{equation}
 \frac{\Gamma\vdash_{\mathcal R}Q_yB(y)\qquad
       \Gamma\vdash_{\mathcal R}B(\eta_B(a))}
      {\Gamma\vdash_{\mathcal R}
         B(\mathsf A(\eta_B(a),n))}
 \quad\mathrm{QE}_{\mathsf A},
 \label{eq:new-qe}
\end{equation}
\begin{equation}
 \frac{\Gamma\vdash_{\mathcal R}D_{x,y}B(x,y)\qquad
       \Gamma,Q_yB(b,y)\vdash_{\mathcal R}F}
      {\Gamma\vdash_{\mathcal R}F}
 \quad\mathrm{DE},
 \label{eq:new-de}
\end{equation}
\begin{equation}
 \frac{\Gamma,B(t,\eta_B(z))\vdash_{\mathcal R}B(t,Sz)}
      {\Gamma\vdash_{\mathcal R}D_{x,y}B(x,y)}
 \quad\mathrm{DI}_{\eta}.
 \label{eq:new-di}
\end{equation}
In \eqref{eq:new-tr}, $p(v),q(v)$ are numerical term templates.  In
\eqref{eq:new-qe}, $a,n$ are arbitrary terms in the conclusion's context.
In DE, $b$ is fresh for $\Gamma,F$ and all fixed parameters of $B$.
In DI, $z$ is fresh for $\Gamma,t$ and all fixed parameters of $B$.
The displayed local assumptions are discharged.  The designated
contradiction is $0=1$.

There are no other primitive axioms, inference rules, induction schemas,
logical connectives, successor axioms, or numerical case principles.
In particular, equality transport in a compound formula is not primitive.

\begin{lemma}[Structural rules and substitution]
\label{lem:new-structural}
Weakening, exchange, contraction, renaming, assumption cut, and capture-free
numerical substitution are admissible in $\mathcal R$.
\end{lemma}
\begin{proof}
Use induction on a finite derivation, renaming eigenparameters before a
substitution or a cut.  A cut replaces each relevant hypothesis leaf by the
supplied derivation, weakened into its local context.  The only new
substitution check is
\[
 (\eta_B(a))\sigma=\eta_{B\sigma}(a\sigma).
\]
It holds because numerical substitution preserves the outermost formula
constructor.  The iterator abbreviation \eqref{eq:new-add} also commutes with
substitution.  Hence QE and DI translate to instances of the same rules;
there is no witness-literal exception to repair.
\end{proof}

\subsection{Equality and the induction bootstrap}
\label{subsec:new-bootstrap}

A useful derived rule is
\begin{equation}
 \frac{\Gamma,B(\eta_B(z))\vdash_{\mathcal R}B(Sz)}
      {\Gamma\vdash_{\mathcal R}Q_yB(y)}
 \quad\mathrm{QI}_{\eta}.
 \label{eq:new-qi}
\end{equation}
To obtain it, apply DI to $D_{x,y}B(y)$ with $x$ unused and witness $0$.
Apply DE to this result; its local hypothesis is precisely $Q_yB(y)$, which
can be exported by Hyp because it is independent of the fresh witness.

\begin{lemma}[Numerical equality and contextual equation induction]
\label{lem:new-ei}
The usual reflexivity, symmetry, transitivity, and numerical congruence rules
are derivable.  For every literal equation $E(y)$,
\begin{equation}
 \frac{\Gamma\vdash_{\mathcal R}E(0)\qquad
       \Gamma,E(z)\vdash_{\mathcal R}E(Sz)}
      {\Gamma\vdash_{\mathcal R}E(t)}
 \quad(z\text{ fresh})
 \label{eq:new-ei}
\end{equation}
is admissible with arbitrary retained side assumptions $\Gamma$.
\end{lemma}
\begin{proof}
Let $J(a)=I(0,a;u.0)$.  I0 gives $J(a)=a$; transport this equation along
itself with template $v=a$ to obtain $a=a$.  The usual equational templates
then give symmetry, transitivity, and numerical congruence.

For induction, the step premise and \eqref{eq:new-qi} give $Q_yE(y)$.
Use \eqref{eq:new-qe} with source $a=0$ and initializer $n=t$.  Its conclusion
is $E(I(0,t;u.Su))$.  I0 and purely equational transport give $E(t)$.
No identity $I(t,0;u.Su)=t$ has been assumed.  Using the counter $t$ and
initializer $0$ in the elimination rule instead would not justify this
bootstrap.
\end{proof}

\begin{lemma}[Elementary iterator identities]
\label{lem:new-addition}
The new calculus proves, with arbitrary parameters,
\begin{equation}
 \begin{aligned}
 \mathsf A(0,n)&=n,&
 \mathsf A(Sa,n)&=S\mathsf A(a,n),\\
 \mathsf A(a,0)&=a,&
 \mathsf A(a,Sn)&=S\mathsf A(a,n),\\
 \mathsf A(a,n)&=\mathsf A(n,a),&
 \mathsf A(a,1)&=Sa.
 \end{aligned}
 \label{eq:new-add-identities}
\end{equation}
\end{lemma}
\begin{proof}
The first row consists of computation instances.  Induct on $a$ for
$\mathsf A(a,0)=a$.  Induct on $a$ again for the initializer-successor
identity: its base is $Sn=Sn$, and its step follows by applying $S$ to the
induction hypothesis.  For commutativity, induct on $a$.  The base is
$n=\mathsf A(n,0)$; the step is
\[
 \mathsf A(Sa,n)=S\mathsf A(a,n)
               =S\mathsf A(n,a)=\mathsf A(n,Sa).
\]
Finally $\mathsf A(a,1)=S\mathsf A(a,0)=Sa$.
All inductions are instances of \eqref{eq:new-ei} on literal equations.
\end{proof}

An atomic closure therefore has the derived one-step elimination
\begin{equation}
 Q_yE(y),\ E(a)\ \vdash_{\mathcal R}\ E(Sa).
 \label{eq:new-step}
\end{equation}
Take initializer $1$ in QE and use \eqref{eq:new-add-identities}.
For a compound matrix, QE at source $0$ and initializer $t$ has the ordinary
conditional-closure effect $Q_yB,B(0)\vdash B(t)$, once transport through
$B$ has been derived.

\begin{lemma}[Full transport and homogeneous binder congruence]
\label{lem:new-transport}
For every formula template $F(v)$,
\begin{equation}
 s=t,\ F(s)\ \vdash_{\mathcal R}\ F(t).
 \label{eq:new-full-tr}
\end{equation}
Moreover, if $B$ and $C$ are uniformly interderivable and have the same root
class (both equations or both compound), their $Q$ closures are
interderivable, and their $D$ closures are interderivable.
\end{lemma}
\begin{proof}
For transport, induct on $F$.  Equations use the primitive rule.  For
$Q_yE(v,y)$ with atomic $E$, introduce the target closure by QI.
Reverse-transport its local $E(t,z)$ to $E(s,z)$, use
\eqref{eq:new-step} for the source closure, and transport the resulting
$E(s,Sz)$ to $E(t,Sz)$.

For $Q_yB(v,y)$ with compound $B$, introduce the target closure from
$B(t,0)$.  The induction hypothesis converts this to $B(s,0)$.
QE for the source closure at initializer $Sz$ gives
$B(s,I(0,Sz;u.Su))$.  The induction hypothesis applied inside the strictly
smaller matrix $B$ removes the I0 redex and changes the parameter $s$ to $t$.
This proves the QI target $B(t,Sz)$.

For $D_{x,y}B(v,x,y)$, eliminate the source binder at a fresh witness $b$.
Introduce the target at that witness.  Convert its local assumption
$B(t,b,\eta_B(z))$ to $B(s,b,\eta_B(z))$ by the induction hypothesis.
At an atomic matrix use \eqref{eq:new-step}; at a compound matrix use QE
with source $0$, initializer $Sz$, and I0.  Transport the successor result
back to the parameter $t$.  Close DI and DE.  Every invoked formula
transport concerns the strictly smaller matrix, so this is not circular.

For homogeneous congruence, use the same constructions, replacing numerical
transport by the two given conversions between $B$ and $C$.  Their root
classes agree, so their local DI assumptions use the same value of $\eta$.
For $D$, eliminate the source at a fresh witness and introduce the target
at that witness.  The argument works uniformly under arbitrary additional
retained assumptions.
\end{proof}

The root-class qualification is substantive.  The new $Q$ is not claimed to
preserve arbitrary equivalence when one matrix is an equation and the other
is compound.  A progressive assertion and a conditional-universal assertion
are different assertions even when their matrices are pointwise equivalent.
The translations below respect this boundary explicitly.

\subsection{Ordinary soundness}
\label{subsec:new-soundness}

Use the pinned repository's scalar theory $T_{\rm eq}$: ordinary
intuitionistic first-order logic with actual equality, the two iterator
computations, numerical coverage
\[
 \forall x\,(x=0\lor\exists y\,x=Sy),
\]
separation $(0=1)\to\bot$, and contextual induction for literal equations.
The repository proves its whole-domain bi-interpretability with iPRA.\cite{bridge}

Define the new ordinary reading recursively by
\begin{equation}
 \begin{aligned}
 \llbracket s=t\rrbracket&=(s=t),\\
 \llbracket Q_yE(y)\rrbracket
   &=\forall y\,(E(y)\to E(Sy))
       &&(E\text{ a literal equation}),\\
 \llbracket Q_yB(y)\rrbracket
   &=\llbracket B(0)\rrbracket\to
                  \forall y\,\llbracket B(y)\rrbracket
       &&(B\text{ compound}),\\
 \llbracket D_{x,y}B\rrbracket
   &=\exists x\,\llbracket Q_yB\rrbracket.
 \end{aligned}
 \label{eq:new-reading}
\end{equation}
This translation preserves all free numerical variables and commutes with
numerical substitution.  It is a finite syntactic translation, not a truth
predicate on numerical formula codes.

\begin{theorem}[Soundness]
\label{thm:new-soundness}
If $\Gamma\vdash_{\mathcal R}F$, then
$\llbracket\Gamma\rrbracket\vdash_{T_{\rm eq}}\llbracket F\rrbracket$.
\end{theorem}
\begin{proof}
Hyp, computations, equational transport, and DE have their ordinary meanings.
Atomic DI is universal introduction of progress, followed by existential
introduction at its witness.  Compound DI is conditional universal
introduction using numerical coverage, followed by existential introduction.

For atomic QE, retain
$H=\forall y(E(y)\to E(Sy))$ and $E(a)$.
Equation induction on $n$ gives $E(\mathsf A(n,a))$: the base is $E(a)$,
and the step uses $H$ at $\mathsf A(n,a)$ and IS.  The pure equational
proofs in Lemma~\ref{lem:new-addition} also hold in $T_{\rm eq}$, so
commutativity gives $E(\mathsf A(a,n))$.  For compound QE the minor is
$\llbracket B(0)\rrbracket$; its conditional-universal major directly
supplies the target instance.  Thus every rule translates soundly.
\end{proof}

\subsection{A precise comparison with the existing native calculus}
\label{subsec:new-native-comparison}

Let $\mathcal N$ denote exactly the current native Eq/$Q$/$D$ calculus
\texttt{NPrf}.  Its $Q_yB$ always means $B(0)\to\forall yB(y)$; its direct
progress rule introduces $Q_yE$ only for equations.  Its terms and numerical
equality are identical to those of $\mathcal R$.

For $E\equiv(s=t)$ define
\begin{equation}
 \begin{aligned}
 C(v,a,b)&:=I(v,a;u.b),\\
 G_vE&:=\bigl(C(v,0,s)=C(v,0,t)\bigr),\\
 L(E)&:=D_{r,v}G_vE,\\
 P(F)&:=\begin{cases}L(F),&F\text{ an equation},\\F,&F\text{ compound}.
              \end{cases}
 \end{aligned}
 \label{eq:new-padding}
\end{equation}
All newly displayed binders are fresh; the update in $C$ ignores $u$, and
$r$ is unused in $L(E)$.

\begin{lemma}[Padding]
\label{lem:new-padding}
In $\mathcal R$, $L(E)$ and $E$ are interderivable under arbitrary contexts.
Consequently $P(F)$ is compound and interderivable with $F$.
\end{lemma}
\begin{proof}
Computations and equality prove $G_0E$ and
$G_{Sz}E\dashv\vdash E$.  From $E$, introduce $L(E)$ by atomic DI, ignoring
its local induction assumption and using $E$ for the successor target.
Conversely, DE on $L(E)$ gives $Q_vG_vE$.  Supply the theorem $G_0E$ and
apply QE with source $0$, initializer $1$.  I0 gives $G_1E$, which reduces
to $E$.  This proof uses no ordinary logical macro.
\end{proof}

Define a translation $T:\mathcal N\to\mathcal R$ by
\begin{equation}
 \begin{aligned}
 T(s=t)&=(s=t),\\
 T(Q_yB)&=Q_yP(TB),\\
 T(D_{x,y}B)&=D_{x,y}P(TB).
 \end{aligned}
 \label{eq:new-T}
\end{equation}
It fixes terms, commutes with numerical substitution, and places a compound
matrix immediately below every translated binder.

\begin{theorem}[Translation into the new calculus]
\label{thm:new-T}
Every $\mathcal N$ derivation translates effectively:
\[
 \Gamma\vdash_{\mathcal N}A
 \quad\Longrightarrow\quad
 T\Gamma\vdash_{\mathcal R}TA.
\]
\end{theorem}
\begin{proof}
Induct on the old derivation.  Hyp, computations, and equational transport
translate literally.  Old QE becomes compound new QE: pad the minor,
specialize at the old target using source $0$, remove the I0 redex, and
unpad.  In DE the opened new assumption is exactly
$Q_yP(TB)(b,y)$, the translation of the old opened $Q$ assumption.
For case DI, unpad the local base, apply the translated premise, and pad
the successor conclusion; the new matrix is compound.

For old direct Q-progress at $E$, the translated target is $Q_yL(E(y))$.
Introduce it by compound QI.  Its local base $L(E(0))$ supplies $E(0)$.
Apply the recovered contextual equation induction to the translated step,
then pad $E(Sz)$ to $L(E(Sz))$.  This discharges the only nonliteral rule
translation.  Each construction is finite and preserves eigenvariable
conditions after fresh renaming.
\end{proof}

\subsection{Logical macros and the order of the completeness proof}
\label{subsec:new-image-logic}

The old ordinary logical macros will be denoted
$\bot_{\mathcal N}$, $\to_{\mathcal N}$, $\wedge_{\mathcal N}$,
$\vee_{\mathcal N}$, $\exists_{\mathcal N}$, and $\forall_{\mathcal N}$.
For specificity they are the macros of the pinned
\texttt{NativeLogic.lean}:\cite{native}
\begin{equation}
 \begin{aligned}
 \bot_{\mathcal N}&=(0=1),\\
 A\to_{\mathcal N}B&=Q_y\operatorname{sel}(y,A,B),\\
 A\wedge_{\mathcal N}B
   &=Q_y\operatorname{gate}
           (y,\operatorname{sel}(\operatorname{par}(y),A,B)),\\
 \exists_{\mathcal N}x\,A(x)&=D_{x,y}\operatorname{gate}(y,A(x)),\\
 A\vee_{\mathcal N}B&=\exists_{\mathcal N}v\,
                                   \operatorname{sel}(v,A,B),\\
 \forall_{\mathcal N}x\,A(x)&=A(0)\wedge_{\mathcal N}Q_xA(x).
 \end{aligned}
 \label{eq:new-old-macros}
\end{equation}
Here $\operatorname{par}(y)=I(y,0;u.C(u,1,0))$.
The old gate acts on equations as $G_v$ in \eqref{eq:new-padding}.
The exact finite recursive definitions of the gate and selector, including
their behavior under $Q$ and $D$, are stated in Appendix~\ref{app:old-macros}.
The selector has zero endpoint $A$ and successor endpoint $B$.
All these operations concern the \emph{old} uniform conditional calculus.
The pinned repository supplies their capture-free definitions and derived
rules.

\begin{lemma}[Logic on image formulas, with arbitrary side context]
\label{lem:new-image-logic}
For old formulas $A,B$, the formulas
$T(A\to_{\mathcal N}B)$, $T(A\wedge_{\mathcal N}B)$, and the analogous
translated macros have their ordinary natural-deduction rules on principal
formulas $TA,TB$.  These derived rules permit an arbitrary retained
$\mathcal R$ context and arbitrary supplied $\mathcal R$ premise proofs.
Image existential elimination also permits an arbitrary new conclusion,
not necessarily an image formula.
\end{lemma}
\begin{proof}
Translate the finite \emph{derived-rule constructions}, rather than merely
closed old theorems.  Their nonlogical steps are replaced by the local
simulations in Theorem~\ref{thm:new-T}.  Their open premise leaves are filled
with the supplied new proofs; assumption cut and weakening
(Lemma~\ref{lem:new-structural}) retain any additional new context.  The
construction never examines the syntax of passive side assumptions.
Thus this is a uniform rule simulation, not an inference from
closed-theorem preservation alone.

For the important existential exception, the old elimination construction
opens its defining $D$ by DE, obtains the witness instance from the opened
gated closure, and cuts this instance into the supplied minor derivation.
It neither analyzes nor introduces the minor's conclusion.  The translated
construction consequently works for every new conclusion $F$, using the
same freshness condition.  This is precisely the structure of the pinned
\texttt{NPrf.exE} proof.
\end{proof}

We do \emph{not} yet assert that the old selectors or logical macros can be
applied unchanged to arbitrary mixed formulas.  Such an assertion would
incorrectly ignore the root distinction in $Q$.  The next normalization
argument first proves that every mixed formula is equivalent to an image.

\subsection{Both formula composites for the entire language}
\label{subsec:new-composites}

Define $J:\mathcal R\to\mathcal N$, fixing all terms and equations, by
\begin{equation}
 \begin{aligned}
 J(Q_yE)&=\forall_{\mathcal N}y\,
                (E(y)\to_{\mathcal N}E(Sy)),
                    &&E\text{ an equation},\\
 J(Q_yB)&=Q_yJB,     &&B\text{ compound},\\
 J(D_{x,y}E)&=\exists_{\mathcal N}x\,
             \forall_{\mathcal N}y\,
                (E(x,y)\to_{\mathcal N}E(x,Sy)),
                    &&E\text{ an equation},\\
 J(D_{x,y}B)&=D_{x,y}JB,&&B\text{ compound}.
 \end{aligned}
 \label{eq:new-J}
\end{equation}
This is a fixed finite syntax transformation.  It maps compound formulas to
compound formulas, preserves free variables, and commutes with numerical
substitution.  Its ordinary reading in $T_{\rm eq}$ is provably equivalent to
\eqref{eq:new-reading}, by the established meanings of the old logical macros.

\begin{lemma}[Atomic closure normalization]
\label{lem:new-atomic-normalization}
For every equation $E(x,y)$, under arbitrary retained contexts,
\begin{equation}
 \begin{aligned}
 Q_yE(x,y)&\dashv\vdash_{\mathcal R}
       T\bigl(\forall_{\mathcal N}y
           (E(x,y)\to_{\mathcal N}E(x,Sy))\bigr),\\
 D_{x,y}E(x,y)&\dashv\vdash_{\mathcal R}
       T\bigl(\exists_{\mathcal N}x\,
             \forall_{\mathcal N}y
           (E(x,y)\to_{\mathcal N}E(x,Sy))\bigr).
 \end{aligned}
 \label{eq:new-atomic-normalization}
\end{equation}
\end{lemma}
\begin{proof}
Write $H(x)$ for the image formula on the first line's right.
From $Q_yE(x,y)$ and the temporary assumption $E(x,z)$, derived one-step
elimination gives $E(x,Sz)$.  The image implication and universal
introductions of Lemma~\ref{lem:new-image-logic} give $H(x)$.  Their
principal component formulas are just equations, which are already fixed
by $T$; no full mixed-language logical rule has been assumed.
Conversely, under $H(x)$ and $E(x,z)$, image universal specialization and
implication elimination give $E(x,Sz)$; atomic QI gives $Q_yE(x,y)$.

For $D$, eliminate the left-hand binder at fresh $b$, normalize its local
$Q_yE(b,y)$ to $H(b)$ by the first part, introduce the image existential
at $b$, and close DE.  In the reverse direction, image existential
elimination supplies a fresh $b$ and $H(b)$.  Under $E(b,z)$ its image
universal and implication rules give $E(b,Sz)$, so atomic DI at witness
$b$ gives $D_{x,y}E$.  Close image existential elimination.  Its conclusion
is allowed to be this arbitrary new formula by the last assertion of
Lemma~\ref{lem:new-image-logic}.
\end{proof}

\begin{theorem}[Exact native-language round trips]
\label{thm:new-roundtrips}
For every old formula $A$ and every new formula $F$,
\begin{equation}
 A\dashv\vdash_{\mathcal N}JT(A),\qquad
 F\dashv\vdash_{\mathcal R}TJ(F).
 \label{eq:new-roundtrips}
\end{equation}
These equivalences are uniform under arbitrary additional retained contexts.
\end{theorem}
\begin{proof}
For the old-side composite, note first that $JL(E)$ is old-equivalent to
$E$.  It asserts existence of an unused witness and progressivity of the
guarded equation $G_vE$.  From $E$, every positive guard instance follows,
so the progress implications hold.  Conversely $G_0E$ is a theorem and
progress at $0$ yields $G_1E$, hence $E$.  This uses ordinary logical
macros in $\mathcal N$, whose correctness is already established.
Now induct on the old formula.  Every binder created by $T$ has a compound
matrix, so $J$ gives its old conditional interpretation.  The preceding
padding equivalence and the induction hypothesis apply under the old
$Q$ and $D$ congruence rules.

For the new-side composite, induct on $F$.  Equations are unchanged.
Atomic $Q$ and $D$ cases are exactly
Lemma~\ref{lem:new-atomic-normalization}.  Suppose the matrix $B$ is
compound.  The induction hypothesis gives $B\dashv\vdash TJB$.
Both $B$ and $TJB$ are compound, since $J$ and $T$ preserve compound roots.
Homogeneous congruence (Lemma~\ref{lem:new-transport}) therefore gives
\[
 Q B\dashv\vdash Q(TJB),\qquad
 D B\dashv\vdash D(TJB).
\]
Because $TJB$ is compound, no padding is inserted on these right sides;
they are exactly the required $TJ$ composites.  There is no appeal to
congruence across unlike root classes.  The induction proves normalization
for every formula before ordinary logic is extended to arbitrary new
formulas, avoiding a circular completeness argument.
\end{proof}

\paragraph{All ordinary logical operations in the new language.}
Now, and only now, define for arbitrary new formulas
\begin{equation}
 \begin{aligned}
 \bot_{\mathcal R}&=(0=1),\\
 F\mathbin{\diamond_{\mathcal R}}G
   &:=T\bigl(JF\mathbin{\diamond_{\mathcal N}}JG\bigr)
       &&(\diamond\in\{\to,\wedge,\vee\}),\\
 \exists_{\mathcal R}x\,F&:=T(\exists_{\mathcal N}x\,JF),\\
 \forall_{\mathcal R}x\,F&:=T(\forall_{\mathcal N}x\,JF).
 \end{aligned}
 \label{eq:new-full-macros}
\end{equation}
The round trips convert arbitrary new principal formulas to their images;
Lemma~\ref{lem:new-image-logic} then supplies the ordinary rules, and the
round trips convert their conclusions back.  For disjunction or existential
elimination with a new conclusion, normalize that conclusion first if
necessary.  All eigenvariable conditions are preserved because $J,T$
preserve free numerical variables.

This gives intuitionistic natural deduction for the entire new language,
including explosion, ordinary numerical coverage, and the quantified
rules.  It adds no primitive logical axioms.  In particular, the system does
not use excluded middle for arbitrary quantified formulas.

\subsection{Bi-interpretability, consistency, and the inherited arithmetic}
\label{subsec:new-final-bridge}

\begin{theorem}[Full consequence equivalence and bi-interpretability]
\label{thm:new-biinterpretability}
The translations $T,J$ preserve derivations in both directions and have the
round trips \eqref{eq:new-roundtrips}.  The ordinary realization of
$\mathcal R$ is equiconsistent and bi-interpretable with full first-order
iPRA.  Both numerical domains are whole, numerical equality is actual
equality, and the numerical composite isomorphisms have graph $x=y$.
\end{theorem}
\begin{proof}
Preservation by $T$ is Theorem~\ref{thm:new-T}.  For $J$, translate a new
derivation into $T_{\rm eq}$ using Theorem~\ref{thm:new-soundness}; its
formulas are provably equivalent to the ordinary readings of their $J$
translations.  The pinned
native completeness and native formula round-trip theorems translate this
ordinary derivation back into $\mathcal N$ and identify each resulting
formula with its $J$ translation.  Equivalently, one may translate each new
rule directly using the old ordinary macros and equation induction; the
only added arithmetic calculation is the QE soundness proof above.
The two exact formula composites are Theorem~\ref{thm:new-roundtrips}.

The remaining scalar-to-iPRA arithmetic is inherited, not reproved by the
new elimination rule.  The repository's chain constructs a scalar
predecessor, a term reflecting actual equality, quantifier-free induction,
pairing, and a compiler for each fixed primitive-recursive description.
The reverse interpretation expresses a fixed scalar iterator by primitive
recursion.  Equation induction proves the two numerical compilation round
trips.  These are the contents of the pinned \texttt{biInterp},
\texttt{native\_iff\_pra}, \texttt{pra\_iff\_native},
\texttt{native\_roundTrip}, and \texttt{pra\_native\_roundTrip} results.\cite{native,bridge}

Composing that existing bi-interpretation with $T,J$ proves the assertion.
Our two translations leave every scalar term and the numerical domain
unchanged, so they add neither a cut nor a quotient to the existing
identity graphs.  The designated contradiction $0=1$ is fixed.  Effective
finite proof translations therefore give equiconsistency.  This is an
external equivalence theorem, not an internal assertion of iPRA's own
consistency.
\end{proof}

\paragraph{What changes compared with the earlier progressive proposal.}
The repository already contains a presentation with a progressive binder
on every matrix, separate one-step and equation-iteration eliminations,
counter-aware primitive recursion, and an explicit Boolean-cases rule.\cite{progressive}
The progressive reading itself is therefore inherited motivation.  The
present contribution is the root-dependent complete reading with the scalar
iterator unchanged, the single reversed-orbit specialization
\eqref{eq:new-qe}, the removal of a separate case principle, and both
round trips for the entire Eq/$Q$/$D$ formula language.  The mixed reading
without the stronger elimination would lose information at atomic closures;
its soundness alone would not establish the theorem above.

\paragraph{Verification status.}
The new metatheorems in this section are written proofs.  They are not claimed
to have been checked in Lean or Rocq.  The existing Lean theorem names identify
the precise inherited bridge.  Executable certificates and checker tests
support concrete implementations; they do not replace the universal proof
or establish a global lower bound on presentation size.

% END INCLUDE

% BEGIN sections/comparison.tex
\section{Complete checkers and measured simplicity}
\label{sec:comparison}

\subsection{The main typed checker and matched controls}

The recommended implementation of $\mathcal R$ is
\path{code/orbit_checker.py}. It reads exactly a JSON object with
\texttt{goal} and \texttt{proof} fields, verifies that the goal is a
closed formula, synthesizes the proof's conclusion, and requires exact
agreement. Term and formula tags, arities, integer indices, binder scopes,
context membership, and eigenvariable conditions are checked. A successful
run exits normally and prints exactly \texttt{ACCEPT} followed by a
newline. Any exception or resource failure rejects the input.

The two comparison modules implement the pinned eight-rule $\mathcal N$
and seven-rule $W$ calculi with the same syntax, substitution operations,
validation, and file interface. They are newly written matched controls;
they are not represented as Python files published upstream. All helpers
and the JSON interface are included in each standalone module.

\begin{center}
\small
\begin{tabular}{@{}lrrrr@{}}
\toprule
Complete typed checker & Tokens & AST & Bytes & Gzip 9 \\
\midrule
Adopted $\mathcal N$, matched & 1,235 & 1,407 & 4,692 & 1,199 \\
New orbit $\mathcal R$ & 1,206 & 1,366 & 4,553 & 1,211 \\
Known $W$, matched & 1,173 & 1,334 & 4,456 & 1,198 \\
\bottomrule
\end{tabular}
\end{center}

These numbers use CPython 3.12.14 and zlib 1.3.2. Genuine docstrings are
removed from the parsed AST before \texttt{ast.unparse} normalization;
comments and formatting are consequently absent. Bytes include the final
newline. Gzip uses compression level 9 with \texttt{mtime=0} and includes
its wrapper. The evidence also records raw source bytes, raw gzip sizes,
and both source hashes. Gzip is informational and is not minimized.

Relative to the adopted eight-rule control, the new checker saves
139 normalized bytes (about 2.96\%), 41 AST nodes, and 29 significant
tokens. It nevertheless costs 97 bytes, 32 AST nodes, and 33 tokens more
than the matched $W$ checker. Its normalized gzip is also larger than
both controls. The uniform introduction and direct substitution theorem
therefore do not establish a new minimum under the user's implementation
metric.

\subsection{A separate compact implementation improvement}

The archive also contains three compact implementations of the
\emph{earlier guarded-source calculus}. They are not implementations of
the new orbit calculus. Their comparison contract is the same as that of
the authenticated incoming guarded-source package: canonical guarded
public statements, with broader administrative proof data. The compact
programs do not themselves reject every noncanonical public statement.
The main typed checker, in contrast, validates its complete Eq/$Q$/$D$
public grammar. These are different input contracts.

\begin{center}
\small
\begin{tabular}{@{}lrrrr@{}}
\toprule
Compact complete module & Tokens & AST & Bytes & Gzip 9 \\
\midrule
Previous token profile & 257 & 285 & 568 & 338 \\
New token profile & \textbf{256} & 283 & 593 & 346 \\
Previous AST profile & 261 & 279 & 603 & 345 \\
New AST profile & 260 & \textbf{278} & 602 & 344 \\
Previous byte profile & 264 & 301 & 522 & 328 \\
New byte profile & 263 & 300 & \textbf{521} & 327 \\
\bottomrule
\end{tabular}
\end{center}

Thus the inspected attained coordinate envelope changes from
\[
 (257\text{ tokens},279\text{ AST nodes},522\text{ bytes})
 \quad\hbox{to}\quad
 (256,278,521).
\]
Each triple is attained by three different programs. No single module
attains all three coordinates. The token profile actually gets longer
in bytes and gzip, which is why reporting only the coordinate envelope
would conceal relevant information.

The byte and AST changes replace the initial administrative truth
\texttt{[]} by raw integer \texttt{0}, and change neutral major fields
in emitted certificates accordingly. These are both spellings of a
true administrative resource, not new numerical axioms. At nonneutral
major premises the original rule tests are unchanged. At a neutral
major, both representations select the same computation or introduction
branch and contribute no mathematical assumption. The old canonical
proof-production construction therefore compiles with the new neutral
spelling, while soundness uses the same interpretation of retained
resources. No complete equivalence claim is made about every malformed
or noncanonical raw byte stream.

The token profile instead turns the recursive continuation into a
\texttt{while} loop and uses an immutable tuple context. Its recursive
major-premise call is retained. The loop state is exactly the old
continuation state: the current goal, retained context, and global fresh
index. Induction on consumed packets identifies the two executions.
The global certificate convention is unchanged.

These operations compose techniques already present in the incoming
packages. They do not constitute a new arithmetic basis or a new logical
axiomatization. Their strict implementation gain is real but very small.
The frozen predecessors, their hashes, the changed certificates, and the
entire new programs are supplied for inspection.

\subsection{What the executable evidence checks}

The recorded reproduction run has 270 process expectations: 122 expected
acceptances and 148 expected rejections, each checked in ordinary or
optimized Python as specified in the manifest. Two additional tests
exercise open sequents with displayed assumptions, and four controls
exercise capture-avoiding substitution. There are zero mismatches.

Positive examples cover both computation equations, reflexivity derived
from computation and transport, equation induction whose hypothesis is
actually used, zero- and nonzero-source orbit specializations, compound
case introduction, and closure extraction through an unused witness.
In particular, a complete closed certificate proves
$I(3,0;u.Su)=3$ using the new induction construction.

Negative examples include a false equation, wrong computation endpoints,
an incorrect orbit antecedent, a future witness, an escaping existential
eigenparameter, unbound iterator state, negative hypothesis indices,
noninteger variable indices, attempted compound induction, attempted
primitive compound transport, malformed instructions, and an open public
goal. These address the rules that could otherwise strengthen the theory
or break scoping.

One open test derives $0=1$ from a \emph{false progressive parity
assumption}. It shows that the new closure retains its progressive
meaning; it is not a closed contradiction proof. Another attempts to use
a compound closure progressively and is rejected. The full assumption
lists and their scope are recorded in \path{tests/open_sequents.json}.

Finite checks support the implementation, while the written
metatheorems establish the universal mathematical claim. No number of
these finite controls would by itself prove soundness, completeness,
bi-interpretability, or a minimality theorem.

% END INCLUDE

% BEGIN sections/audit.tex
\section{Proof audit and limits of the reduction}
\label{sec:proof-audit}

The seven-rule presentation changes the meaning of an atomic closure.
Consequently its equivalence with the native kernel needs two formula
translations.  It is not an identity-language deletion of one native rule.
This section records the dependency checks that matter for that claim.

\subsection{The order of the arithmetic bootstrap}

Put $A(a,n)=I(a,n;u.Su)$.  The counter is the first argument.
The following order avoids using a derived rule in its own proof.
\begin{enumerate}
\item Computation and equational transport give numerical reflexivity,
      symmetry, transitivity, and term congruence.  For example, if
      $W(t)=I(0,t;u.0)$, the equation $W(t)=t$ is a computation;
      transporting that equation along itself yields $t=t$.
\item A vacuous witness in DI, followed by DE, gives root-directed
      Q-introduction.  This construction uses no numerical induction.
\item Given an equation-induction step and a base proof $E(0)$,
      Q-introduction gives the progressive closure of $E$.  Orbit
      elimination with source $0$ and initializer $t$ gives
      $E(A(0,t))$.  The computation $A(0,t)=t$ yields $E(t)$.
\item Only after this recovery of contextual equation induction do we
      prove $A(a,1)=Sa$.  Orbit elimination with source $a$ and initializer
      $1$ then gives the immediate progress rule
      $\mathsf Q_yE(y),E(a)\vdash E(Sa)$.
\item Full formula transport is proved by structural induction, using
      this immediate progress rule at atomic closure matrices and the
      ordinary base-to-instance elimination at compound matrices.
\end{enumerate}

In particular, the proof does not begin by assuming $I(t,0;u.Su)=t$.
That identity is a consequence of the recovered induction.  Replacing
the target $I(a,n;u.Su)$ in the new elimination rule by
$I(n,a;u.Su)$ would destroy the direct computation in the third item.
The intended rule and the executable representation must retain the
same argument order.

The addition commutation law used to establish soundness in the ordinary
scalar theory is harmless at that stage: the target scalar theory already
has contextual equation induction.  That soundness argument is separate
from the internal recovery of induction in the new calculus.

\subsection{Why equational transport suffices}

For a transport template headed by $\mathsf Q$, the induction hypothesis
applies only to its strictly smaller matrix.  Reverse the supplied
numerical equality to convert the target introduction hypothesis to the
source hypothesis.  Apply the source closure, and convert the resulting
successor instance forwards.  The atomic case uses the immediate progress
rule established above; the compound case uses a zero source in orbit
elimination.  Root-directed Q-introduction closes the argument.

For a template headed by $\mathsf D$, first eliminate the source
$\mathsf D$ at a fresh witness.  Introduce the target $\mathsf D$ at that
same witness, performing the preceding conversion inside its matrix.
Both eigenparameters are chosen outside the transport equality, the
retained context, and the fixed matrix parameters.  Every invocation of
formula transport in these arguments concerns the smaller matrix.
Thus this proof does not silently replace the primitive equational
transport rule by a stronger logical axiom.

The same root invariance is useful for substitution: substituting numerical
terms never turns a literal equation into a compound formula.  The new
introduction and elimination modes therefore survive numerical substitution
without a special witness-zero case.  Renaming eigenparameters first gives
the usual substitution proof by induction on derivations.

\subsection{Root-sensitive closure is not unrestricted congruence}

Padding gives an equation $E$ and a compound formula $P(E)$ that are
interderivable.  Nevertheless one must not infer
\[
               \mathsf Q_yE(y)\ \dashv\vdash\ 
               \mathsf Q_yP(E(y))
\]
merely from that fact.  The first closure is progressive and the second
is conditional.  Numerical transport is legitimate because it preserves
the matrix root.  The binder-congruence lemma used in the formula
composites requires the same mode on its two sides.

The new-side round trip can be justified without assuming logical
completeness for all new formulas in advance.  First translate the old
logical rules for formulas in the image of the old-to-new translation,
allowing an arbitrary passive new context.  Atomic new closures normalize
to image formulas expressing
$\forall y(E(y)\mathbin{\to}E(Sy))$, since $E(y)$ and $E(Sy)$ are
already image formulas.  Atomic witnessed closures then normalize using
the image existential.  Compound cases use the structural induction
hypothesis and same-mode binder congruence.  The translations of compound
formulas remain compound, so that last use of congruence is valid.

One detail is necessary in the atomic witnessed case: elimination of an
image existential must permit an arbitrary new conclusion, not only a
conclusion already known to be an image formula.  Its translated native
macro elimination has precisely this property.  It opens a $\mathsf D$
resource, removes padding and the gate from the witness assumption, and
uses the supplied minor derivation.  The conclusion is only carried
through DE; it is not inspected or introduced beneath a mode-sensitive
binder.  This prevents a circular appeal to the desired normalization
theorem.

\subsection{A concrete obstruction to keeping the old atomic reading}

The change of reading is essential.  A short counterexample uses the
primitive recursive parity term
\[
 p(n)=I(n,0;u.I(u,1;v.0)),\qquad E(y)\equiv p(y)=1.
\]
Direct computations give $p(0)=0$, $p(1)=1$, and $p(2)=0$ in the
standard natural numbers.  Thus
\[
 E(0)\to\forall y E(y)
 \quad\text{is true, whereas}\quad
 \forall y\bigl(E(y)\to E(Sy)\bigr)
 \quad\text{is false}.
\]
Indeed, the old conditional antecedent is false, while the proposed
progress condition fails at $y=1$.

The new orbit elimination would therefore be unsound if its atomic major
kept the old conditional meaning: take source $1$ and initializer $1$.
Its minor $E(1)$ holds and its target is $E(I(1,1;u.Su))=E(2)$, which
fails.  Likewise, merely exporting an ordinary conditional closure from a
progressive witnessed formula discards information needed for a full
formula round trip.  The stronger atomic elimination and the changed
atomic reading must be adopted together.

\subsection{What the comparison establishes}

The new calculus has unrestricted native formula syntax, a single
root-directed introduction, and a direct substitution argument.  These
are structural improvements over the guarded and witness-sensitive
presentations.  They do not by themselves prove a smaller standalone
Python checker.  The orbit elimination has a richer numerical target,
and a faithful implementation must count that cost.  The report's
checker measurements must therefore distinguish this calculus from any
separate compact implementation of an earlier equivalent presentation.

The arguments above are written metatheoretic proofs.  Finite executable
certificates check particular derivations and implementation behavior;
they are not a new Lean or Rocq verification of either formula
interpretation.  No absolute minimality theorem follows from the
seven-rule count or the measured source sizes.

% END INCLUDE

\appendix
% BEGIN sections/old_macros.tex
\section{Explicit logical macros in the reference language}
\label{app:old-macros}

This appendix fixes the finite macros used by the interpretations. All
operators in this appendix belong to the old conditional language
$\mathcal N$. They must be expanded there before applying the translation
$T$. Applying these definitions unchanged to arbitrary new mixed
closures would be incorrect. These definitions reproduce the interface
of \texttt{NativeGate.lean} and \texttt{NativeLogic.lean}
\cite{native}.

Write $C(v,a,b)=I(v,a;u.b)$, where $u$ is unused, and $1=S0$.
Bindings below are freshened so that no displayed passive parameter is
captured. Define an equation--formula selector $U_v(s,t;B)$ by
\[
\begin{aligned}
U_v(s,t;r=w)&=(C(v,s,r)=C(v,t,w)),\\
U_v(s,t;Q_y B)&=Q_y U_v(C(y,0,s),C(y,0,t);B),\\
U_v(s,t;D_{x,y} B)&=D_{x,y}U_v(C(y,0,s),C(y,0,t);B).
\end{aligned}
\]
Its two endpoints are the equation $s=t$ and the formula $B$.
Put $G_v A=U_v(0,0;A)$. Then $G_0 A$ is provable and
$G_{Sz}A\eqv A$. A selector between arbitrary formulas is
\[
\begin{aligned}
M_v(s=t,B)&=U_v(s,t;B),\\
M_v(Q_y A,B)&=Q_y M_v(A,G_yB),\\
M_v(D_{x,y} A,B)&=D_{x,y}M_v(A,G_yB).
\end{aligned}
\]
This recursion terminates on the depth of its first argument; the
separate recursion defining $U$ terminates on its formula argument.
For every $A,B$, the old calculus derives
$M_0(A,B)\eqv A$ and $M_{Sz}(A,B)\eqv B$.

Let $p(y)=I(y,0;u.C(u,1,0))$. The logical operations are
\[
\begin{aligned}
\bot_*&=(0=1),\\
A\Rightarrow_* B&=Q_y M_y(A,B),\\
A\wedge_* B&=Q_y G_y M_{p(y)}(A,B),\\
\exists_* x\,A(x)&=D_{x,y}G_y A(x),\\
A\vee_*B&=\exists_*v\,M_v(A,B),\\
\forall_*x\,A(x)&=A(0)\wedge_* Q_x A(x).
\end{aligned}
\]
All new bound variables are fresh. These are definitions by finite
syntax, not additional primitive logical axioms or rules.

\paragraph{Why they have the logical rules.}
Implication introduction uses the zero selector as its local hypothesis
and the successor selector as its conclusion; elimination specializes
at $1$. Existential introduction chooses the witness and uses the
successor gate. Existential elimination uses $D$ elimination, the
zero-gate theorem, and specialization at $1$; its final conclusion can
be any formula, independently of the existential matrix.

For conjunction, the gate supplies the base theorem and the computations
$p(1)=1$, $p(2)=0$ recover the two conjuncts. Numerical cases establish
the selector at an arbitrary control from proofs of both alternatives.
Disjunction uses witnesses $0$ and $1$; its elimination uses the derived
case rule on an implication to the desired conclusion. The universal
macro has its zero instance as one conjunct and every instance as the
other conjunct's closure. Finally, $0=1$ and the selector computations
give $s=C(0,s,t)=C(1,s,t)=t$ for arbitrary terms; structural induction
on formulas gives falsity elimination.

The full proofs and their substitution lemmas are preserved in the
pinned source files included with the package. The new theory's
normalization proof uses the translated versions of these derived-rule
constructions on $T$-image principal formulas before it defines the
ordinary logical operations on arbitrary new formulas.

% END INCLUDE

% BEGIN sections/arithmetic_bridge.tex
\section{The arithmetic interface and the inherited bridge}
\label{app:arithmetic-bridge}

The reference theory $\Teq$ is intuitionistic first-order equality logic
over scalar iterator terms, with the two computation equations,
numerical coverage
\[
\forall x\,(x=0\vee\exists y\,x=Sy),
\]
separation $(0=S0)\to\bot$, and induction for literal equations under
side assumptions independent of the induction variable. Its equality
is ordinary equality. The reference iPRA has a symbol for each fixed
primitive recursive definition, its defining equations, successor
separation and injectivity, and quantifier-free induction over ordinary
intuitionistic first-order logic.

The pinned development establishes $\mathcal N\simeq\Teq\simeq\ipra$
\cite{bridge}. The arithmetic part is retained unchanged here. Its
substantive stages are:
\begin{enumerate}
\item Construct a predecessor in the scalar iterator language. Since the
update binds only its previous state, this is a real construction; the
counter variable is not silently made available to the update.
\item Derive equality reflection into a term and Boolean operations.
Replace a quantifier-free formula by its characteristic equation, which
transfers literal-equation induction to quantifier-free induction.
\item Construct pairing and projections, and compile a counter-aware
primitive recursion by iterating a state that carries both the counter
and the accumulated value.
\item Translate each fixed primitive recursive definition and each
scalar iterator term in the opposite direction, proving their defining
equations and both composites by equation induction.
\end{enumerate}

The archive contains the corresponding complete written sources
\path{scalar_predecessor.tex}, \path{equality_reflection.tex}, and
\path{primitive_recursion.tex} in the pinned reference subtree, together
with the ordinary bridge and the native logical basis. These files are
included to expose the mathematical dependency; the new report does
not claim to have discovered these arithmetic constructions.

There is no quantified code for an arbitrary primitive recursive
program and no total universal primitive recursive evaluator in the
theory. A compiler acts externally on one finite definition at a time.
This distinction matters: adjoining a purported total universal
evaluator would change the foundational argument.

Both new translations between $\mathcal R$ and $\mathcal N$ fix every
numerical term and literal equation. Composing them with the reference
arithmetic translations therefore retains whole numerical domains and
actual equality. The graph of each numerical interpretation composite
is provably $x=y$. Formula composites are proved for all formulas, so
the conclusion is bi-interpretability, not merely a correspondence of
closed equations or an equiconsistency argument.

The precise existing Lean endpoint names are
\begin{quote}\small\ttfamily
native\_iff\_pra,\quad pra\_iff\_native,\\
native\_roundTrip,\quad pra\_native\_roundTrip,\\
native\_inconsistent\_iff.
\end{quote}
Their source was retrieved and
authenticated at the pinned commit. Lean/Lake was not available in this
run, so these existing files were inspected rather than rebuilt, and no
new Lean theorem about $\mathcal R$ is claimed.

% END INCLUDE

% BEGIN sections/artifacts.tex
\section{Artifacts and reproducibility}

The source archive contains the following main items.
\begin{center}
\small
\begin{tabular}{@{}p{.45\textwidth}p{.48\textwidth}@{}}
\toprule
Path & Purpose \\
\midrule
\path{ipra_uniform_progress.tex}, \path{ipra_uniform_progress.pdf}
& Self-contained LaTeX article and compiled report. \\
\path{report.tex}, \path{sections/} & Editable modular sources. \\
\path{code/orbit_checker.py} & Complete typed checker for the new theory. \\
\path{code/native_checker.py}, \path{code/witness_checker.py}
& Matched standalone old-kernel controls. \\
\path{code/ancillary/} & Separate compact implementation improvements. \\
\path{tests/}, \path{evidence/} & Certificates, untrusted producer, reproduction, measurements, and exact outcomes. \\
\path{reference/} & Authenticated source files and the inherited arithmetic and logic bridge. \\
\path{build.py}, \path{CLAIMS.json}, \path{SHA256SUMS}
& Article rebuild, exact scope, and package integrity. \\
\bottomrule
\end{tabular}
\end{center}

From the extracted package directory:
\begin{lstlisting}[language=bash]
python3 code/orbit_checker.py < tests/certificates/identity_at_three_orbit.json
python3 -B tests/reproduce.py
python3 build.py
\end{lstlisting}
The first command prints \texttt{ACCEPT}. The second regenerates the main
certificates, performs the documented checks, and rewrites the evidence
with the current Python and compression versions. The last assembles the
standalone source and runs \texttt{pdflatex} enough times to settle
references. Python execution uses only the standard library. PDF generation
requires the standard TeX packages declared in the source.

The complete flattened source can also be compiled by itself:
\begin{lstlisting}[language=bash]
pdflatex -interaction=nonstopmode -halt-on-error ipra_uniform_progress.tex
pdflatex -interaction=nonstopmode -halt-on-error ipra_uniform_progress.tex
\end{lstlisting}

The arithmetic and native/iPRA theorem sources are authenticated by their
Git blob SHA-1 identifiers as well as SHA-256 content hashes.
\path{reference/PROVENANCE.json} records their immutable paths and commit.
\path{evidence/incoming_archives.json} records the six inspected ZIP blob
identities. The source repository uses the MIT No Attribution license,
included as \path{reference/LICENSE-ProveIt.txt}.

\clearpage
\subsection{The seven proof-instruction branches}

The following is the actual rule-synthesis function from the standalone
typed checker. Its syntax validators, capture-avoiding substitution,
scope operations, and file interface are also included in the measured
module; they are not hidden dependencies.
\begin{lstlisting}[language=Python]
def derive(proof, context=(), scope=0):
    require(type(proof) is list and proof)
    rule, *args = proof
    if rule == "hyp":
        index, = args
        require(type(index) is int and 0 <= index < len(context))
        return context[index]
    if rule == "zero":
        seed, body = args
        require(term(seed, scope) and term(body, scope + 1))
        return ["=", ["i", 0, seed, body], seed]
    if rule == "succ":
        counter, seed, body = args
        require(term(counter, scope) and term(seed, scope) and term(body, scope + 1))
        return ["=", ["i", ["s", counter], seed, body], instance(body, ["i", counter, seed, body])]
    if rule == "tr":
        matrix, major, minor = args
        require(formula(matrix, scope + 1) and matrix[0] == "=")
        equality = derive(major, context, scope)
        require(equality[0] == "=")
        require(derive(minor, context, scope) == instance(matrix, equality[1]))
        return instance(matrix, equality[2])
    if rule == "use":
        source, seed, major, minor = args
        require(term(source, scope) and term(seed, scope))
        closure = derive(major, context, scope)
        require(closure[0] == "q")
        matrix = closure[1]
        source = source if matrix[0] == "=" else 0
        require(derive(minor, context, scope) == instance(matrix, source))
        return instance(matrix, ["i", source, seed, ["s", ["v", 0]]])
    if rule == "elim":
        conclusion, major, continuation = args
        require(formula(conclusion, scope))
        paired = derive(major, context, scope)
        require(paired[0] == "d")
        local = (["q", paired[1]],) + tuple(lift(item) for item in context)
        require(derive(continuation, local, scope + 1) == lift(conclusion))
        return conclusion
    if rule == "intro":
        matrix, witness, premise = args
        require(formula(matrix, scope + 2) and term(witness, scope))
        body = substitute(matrix, 1, lift(witness))
        source = ["v", 0] if body[0] == "=" else 0
        local = (fresh_instance(body, source),) + tuple(lift(item) for item in context)
        require(derive(premise, local, scope + 1) == fresh_instance(body, ["s", ["v", 0]]))
        return ["d", matrix]
    raise ValueError("Unknown inference rule")
\end{lstlisting}

% END INCLUDE

% BEGIN sections/references.tex
\begin{thebibliography}{9}
\bibitem{baseline}
Vladimir Reshetnikov and ProveIt contributors.
\emph{Concise presentations of intuitionistic primitive recursive arithmetic}.
Repository snapshot
\path{a4e5e96c456303cf231eab2aeeb3291d8889833f}, inspected November 11, 2026.
\href{https://github.com/VladimirReshetnikov/ProveIt/tree/a4e5e96c456303cf231eab2aeeb3291d8889833f/Logic/PrimitiveRecursiveArithmetic/ConciseIPRA}{ConciseIPRA directory and README}.

\bibitem{native}
ProveIt contributors.
\emph{Native scalar syntax, deduction, gates, and logical macros}.
\href{https://github.com/VladimirReshetnikov/ProveIt/blob/a4e5e96c456303cf231eab2aeeb3291d8889833f/Logic/PrimitiveRecursiveArithmetic/ConciseIPRA/Lean/ConciseIPRA/Native.lean}{Native.lean},
\href{https://github.com/VladimirReshetnikov/ProveIt/blob/a4e5e96c456303cf231eab2aeeb3291d8889833f/Logic/PrimitiveRecursiveArithmetic/ConciseIPRA/Lean/ConciseIPRA/NativeGate.lean}{NativeGate.lean}, and
\href{https://github.com/VladimirReshetnikov/ProveIt/blob/a4e5e96c456303cf231eab2aeeb3291d8889833f/Logic/PrimitiveRecursiveArithmetic/ConciseIPRA/Lean/ConciseIPRA/NativeLogic.lean}{NativeLogic.lean}, same snapshot.

\bibitem{bridge}
ProveIt contributors.
\emph{Native/iPRA bi-interpretability and its arithmetic foundations}.
\href{https://github.com/VladimirReshetnikov/ProveIt/blob/a4e5e96c456303cf231eab2aeeb3291d8889833f/Logic/PrimitiveRecursiveArithmetic/ConciseIPRA/Lean/ConciseIPRA/NativeBiInterpretation.lean}{NativeBiInterpretation.lean} and
\href{https://github.com/VladimirReshetnikov/ProveIt/tree/a4e5e96c456303cf231eab2aeeb3291d8889833f/Logic/PrimitiveRecursiveArithmetic/ConciseIPRA/Research/certified_zero/reference}{certified-zero reference proofs}, same snapshot.

\bibitem{witness}
ProveIt contributors.
\emph{One introduction rule}.
\href{https://github.com/VladimirReshetnikov/ProveIt/blob/a4e5e96c456303cf231eab2aeeb3291d8889833f/Logic/PrimitiveRecursiveArithmetic/ConciseIPRA/Lean/ConciseIPRA/NativeOneIntro.lean}{NativeOneIntro.lean}, same snapshot.

\bibitem{guarded}
ProveIt contributors.
\emph{A guarded core with a deterministic introduction rule}.
\href{https://github.com/VladimirReshetnikov/ProveIt/tree/a4e5e96c456303cf231eab2aeeb3291d8889833f/Logic/PrimitiveRecursiveArithmetic/ConciseIPRA/Research/guarded_core}{Research/guarded\_core}, same snapshot.

\bibitem{source}
ProveIt contributors.
\emph{Source-selected elimination and zero-source admissibility}.
\href{https://github.com/VladimirReshetnikov/ProveIt/tree/a4e5e96c456303cf231eab2aeeb3291d8889833f/Logic/PrimitiveRecursiveArithmetic/ConciseIPRA/Research/source_dispatch}{Research/source\_dispatch}, same snapshot.

\bibitem{progressive}
ProveIt contributors.
\emph{A progressive binder presentation of iPRA}.
\href{https://github.com/VladimirReshetnikov/ProveIt/blob/a4e5e96c456303cf231eab2aeeb3291d8889833f/Logic/PrimitiveRecursiveArithmetic/ConciseIPRA/Research/research/theory/progressive_binder.tex}{Research/research/theory/progressive\_binder.tex}, same snapshot.

\bibitem{incoming}
ProveIt research intake.
Six incoming iPRA archives, November 11, 2026.
\href{https://github.com/VladimirReshetnikov/ProveIt/tree/a4e5e96c456303cf231eab2aeeb3291d8889833f/Logic/PrimitiveRecursiveArithmetic/ConciseIPRA/docs/incoming}{docs/incoming}, same snapshot.
The ZIP Git blob identities are recorded in the accompanying evidence.
\end{thebibliography}

% END INCLUDE

\end{document}
